\documentclass[11pt,a4paper]{article}
\usepackage[T1]{fontenc}
\usepackage[utf8]{inputenc}
\usepackage[french]{babel}
\usepackage{lmodern}
\usepackage[margin=27mm]{geometry}
\usepackage{xcolor}
\definecolor{navy}{RGB}{23,49,71}
\definecolor{muted}{RGB}{78,89,98}
\newcommand{\editionlink}[2]{\leavevmode\pdfstartlink user{/Subtype /Link /A << /S /URI /URI (#1) >>}\textcolor{navy}{#2}\pdfendlink}
\pagestyle{empty}
\begin{document}
\color{navy}
\noindent\textsc{Bibliothèque de géométrie algébrique}\par
\vspace{22mm}
\noindent{\fontsize{32}{38}\selectfont\bfseries Éléments de géométrie\par algébrique}\par
\vspace{9mm}
\noindent{\Large Alexandre Grothendieck et Jean Dieudonné}\par
\vspace{18mm}
\noindent{\LARGE Édition française}\par
\vspace{3mm}
\noindent{\large Texte diplomatique --- lecteur cumulatif EGA I--IV}\par
\vspace{10mm}
\color{muted}
\noindent Cette édition réunit les volumes I, II, III-1, III-2 et IV-1 à IV-4,
y compris les préliminaires (EGA~0) qui leur sont associés.
Les signets et les références internes permettent de parcourir le lecteur.\par
\vspace{5mm}
\noindent Édition diplomatique établie à partir des volumes imprimés
diffusés par NUMDAM.\par
\vspace{9mm}
\noindent\editionlink{https://doi.org/10.5281/zenodo.21921588}{Édition française historique : DOI 10.5281/zenodo.21921588}\par
\vspace{3mm}
\noindent\editionlink{https://doi.org/10.5281/zenodo.21921591}{Édition anglaise historique : DOI 10.5281/zenodo.21921591}\par
\vspace{3mm}
\noindent\editionlink{https://github.com/KokunoYumeto/ega-fr}{Sources, fichiers LaTeX et versions : github.com/KokunoYumeto/ega-fr}\par
\vfill
\small
\noindent Préparation des sources et assemblage : OpenAI Codex,
GPT-6 Astra, effort Ultra.\par
\vspace{3mm}
\noindent Le PDF est accompagné du texte LaTeX et du projet éditable,
avec les fichiers nécessaires à la compilation.\par

% BEGIN RIGHTS AND AI TRANSCRIPTION NOTICE
\clearpage
\begingroup
\small
\noindent\textbf{Droits, disponibilité et provenance}

Les droits sur les œuvres originales et sur toute contribution protégée à ces traductions et transcriptions dérivées restent ceux de leurs titulaires respectifs. Dans la mesure où l'éditeur détient des droits sur la traduction, la transcription, les corrections et la composition produites par IA, il les dédie au domaine public sous CC0 1.0 Universel (https://creativecommons.org/publicdomain/zero/1.0/deed.fr) ; cela n'affecte pas les droits sur les œuvres originales et n'accorde aucune autorisation de les utiliser. Les mentions de droits de tiers et les autorisations déjà accordées restent applicables.

Ces éditions sont mises à disposition gratuitement, à des fins non commerciales et exclusivement pédagogiques, dans la présomption que les titulaires de droits ne s'opposeraient pas à la libre disponibilité de ces traductions et transcriptions aux côtés des originaux français sur NUMDAM (www.numdam.org/journals/PMIHES/). Cette présomption ne signifie pas qu'une autorisation a été obtenue ; le libre accès aux originaux n'autorise pas, à lui seul, leur traduction, leur adaptation ou leur republication.

Si vous détenez des droits sur un élément présenté ici et vous opposez à sa disponibilité, veuillez ouvrir une demande sur GitHub (github.com/KokunoYumeto/ega-fr/issues). Le responsable de la publication retirera l'élément concerné à la demande du titulaire des droits.

La transcription moderne, les corrections, la composition et les notes mathématiques sont produites par IA. Les travaux documentés comprennent ceux d'OpenAI Codex, GPT-6 Astra, effort Ultra. Les indications de provenance des contributions antérieures restent applicables. Les textes originaux restent l'œuvre d'Alexandre Grothendieck et Jean Dieudonné ; ils ne sont pas attribués à l'IA. Cette édition n'est approuvée ni par les auteurs, ni par l'IHÉS, ni par NUMDAM.
\par
\endgroup
\clearpage
% END RIGHTS AND AI TRANSCRIPTION NOTICE
\end{document}
